\documentclass[11pt]{article}
\usepackage[margin=1in]{geometry}
\usepackage{amsmath,amssymb,amsthm}
\usepackage{graphicx}
\usepackage{booktabs}
\usepackage{xcolor}
\usepackage[colorlinks=true,urlcolor=blue,linkcolor=black]{hyperref}
\usepackage{listings}

\lstset{
  basicstyle=\ttfamily\footnotesize,
  breaklines=true,
  columns=fullflexible,
  keepspaces=true,
  frame=single,
  framesep=4pt,
  rulecolor=\color{gray!50},
  showstringspaces=false,
}

\newcommand{\E}{\mathbb{E}}
\newcommand{\VL}{V_L}

\title{\textbf{SPY Volatility Study}\\[2pt]
\large Moving Average, EWMA, and GARCH(1,1)}
\author{Name: Nabizhan Bobokhanov}
\date{\today}

\begin{document}
\maketitle

\begin{center}
\fbox{\parbox{0.92\textwidth}{\centering
\textbf{Code repository:} \url{https://github.com/nabizhka/spy-volatility-hw}\\
\small (all code is also listed in Appendix~A.)}}
\end{center}
\bigskip

\noindent\textbf{Data.} Daily SPY closing prices, adjusted for dividends, over
2021-08-01 to 2026-07-30, pulled with \texttt{yfinance} (Stooq fallback) and
cached to \texttt{spy\_prices.csv}. Daily log returns
$r_t=\ln(P_t/P_{t-1})$. Annualization uses $\Delta=1/252$, so
$\sigma_{\text{ann}}=\sigma_{\text{day}}/\sqrt{\Delta}=\sqrt{252}\,\sigma_{\text{day}}$.

\section{Moving-average and EWMA annualized volatility}
\textbf{Method.} The moving-average (MA) estimator averages the last
$n=100$ squared returns; the EWMA estimator applies exponential weights with
$\lambda=0.94$, initialised at the sample variance of the first 100 returns
(zero-mean convention, following RiskMetrics/Hull):
\[
v^{\text{MA}}_t=\frac1n\sum_{i=0}^{n-1}r_{t-i}^2,
\qquad
v^{\text{EWMA}}_t=\lambda v^{\text{EWMA}}_{t-1}+(1-\lambda)r_t^2,
\qquad
\sigma^{\text{ann}}_t=\sqrt{v_t/\Delta}.
\]

\textbf{Pseudocode.}
\begin{lstlisting}[language=Python]
r <- log_returns(SPY_close)                 # r_t = ln(P_t / P_{t-1})
# (a) MA(100)
for t with >= n prior returns:
    v_MA[t]   <- mean(r[t-n+1 .. t]^2)
    vol_MA[t] <- sqrt(v_MA[t] / Delta)
# (b) EWMA(0.94)
v <- sample_var(r[0..99])                    # initial variance
for t = 100 .. N-1:
    v <- lambda*v + (1-lambda)*r[t]^2
    vol_EWMA[t] <- sqrt(v / Delta)
plot vol_MA, vol_EWMA on one figure
\end{lstlisting}

\begin{figure}[h]
\centering
\includegraphics[width=\textwidth]{figures/task1_ma_ewma.png}
\caption{Annualized volatility of SPY daily log returns: MA(100) vs
EWMA($\lambda=0.94$).\; \textcolor{red}{\small [Figure regenerates from real
SPY data when you run \texttt{python src/main.py}; the version shown here was
produced by the test harness.]}}
\end{figure}

\textbf{Discussion.} Both estimators track the same volatility regimes, but the
MA(100) is smooth and slow: a large return enters and leaves the 100-day window
with equal weight, producing the characteristic ``boxcar'' plateaus and a lag
of up to 100 days. The EWMA reacts immediately to new returns and decays them
geometrically (weight $(1-\lambda)\lambda^{k}$), so it spikes higher and falls
back faster around turbulent episodes. The EWMA effective memory is
$\approx 1/(1-\lambda)\approx 17$ days, far shorter than the MA window, which is
why it is noisier but more current.

\section{Fitting GARCH(1,1)}
\textbf{Method.} On the window 2025-06-01 to 2026-05-31 we fit, by maximum
likelihood, a constant-mean GARCH(1,1) with Normal innovations:
\[
r_t=\mu+\varepsilon_t,\quad \varepsilon_t=\sigma_t z_t,\ z_t\sim N(0,1),\qquad
\sigma_t^2=\omega+\alpha\,\varepsilon_{t-1}^2+\beta\,\sigma_{t-1}^2 .
\]
Returns are scaled to percent ($\times100$) for numerical conditioning; the
fitted $\omega$ is rescaled back by $100^2$. The long-run daily variance is
$\VL=\omega/(1-\alpha-\beta)$.

\textbf{Pseudocode.}
\begin{lstlisting}[language=Python]
r_fit <- r over 2025-06-01 .. 2026-05-31
x <- 100 * r_fit
(mu, omega, alpha, beta) <- MLE fit of GARCH(1,1) to x   # arch_model(...).fit()
persistence <- alpha + beta
VL_daily    <- (omega/100^2) / (1 - alpha - beta)
\end{lstlisting}

\textbf{Result.} Running \texttt{src/main.py} prints the full \texttt{arch}
summary and the table below. \emph{Replace the example values with the numbers
from your own data run.}
\begin{center}
\begin{tabular}{lc}
\toprule
Quantity & Fitted value \\
\midrule
$\omega$ (per day, decimal var) & \textit{fill from run} \\
$\alpha$                        & \textit{fill from run} \\
$\beta$                         & \textit{fill from run} \\
persistence $\alpha+\beta$      & \textit{fill from run} \\
long-run daily vol $\sqrt{\VL}$ & \textit{fill from run} \\
long-run annual vol             & \textit{fill from run} \\
\bottomrule
\end{tabular}
\end{center}
For a typical broad-equity index one expects high persistence
$\alpha+\beta$ close to (but below) 1, a small $\alpha$ and a large $\beta$, and
a long-run annualized vol in the mid-teens.

\section{Monte-Carlo forecast vs realized estimates}
\textbf{Method.} From the fitted model we simulate the variance process forward
over the $H$ trading days in 2026-06-01 to 2026-07-30. Each of $M=100$
simulations draws $z_t\sim N(0,1)$, sets $\varepsilon_t=\sigma_t z_t$, and
iterates $\sigma_{t+1}^2=\omega+\alpha\varepsilon_t^2+\beta\sigma_t^2$, seeded by
the last in-sample conditional variance. We overlay the analytic expected path
$\E[\sigma_{t+k}^2]=\VL+(\alpha+\beta)^k(\sigma_t^2-\VL)$ and the realized
MA(100) and EWMA(0.94) estimates over the same dates.

\textbf{Pseudocode.}
\begin{lstlisting}[language=Python]
H <- trading days in 2026-06-01 .. 2026-07-30
s2_0 <- last in-sample conditional variance
for m = 1 .. M:                      # M = 100 paths
    s2 <- s2_0
    for k = 1 .. H:
        z <- N(0,1);  eps <- sqrt(s2)*z
        vol[m,k] <- sqrt(s2 / Delta)          # annualized
        s2 <- omega + alpha*eps^2 + beta*s2   # GARCH step forward
realized_MA, realized_EWMA over same window
plot all paths, sim mean, analytic E[vol], realized MA & EWMA
\end{lstlisting}

\begin{figure}[h]
\centering
\includegraphics[width=\textwidth]{figures/task3_forecast_vs_realized.png}
\caption{GARCH(1,1) $M{=}100$ simulated annualized-volatility paths (grey),
their mean (red), the analytic expected forecast (dashed), and realized MA /
EWMA over 2026-06 to 2026-07.\; \textcolor{red}{\small [Regenerates from real
SPY data on \texttt{python src/main.py}.]}}
\end{figure}

\textbf{Discussion.} The individual simulated paths fan out around a
\emph{mean} that decays smoothly toward the long-run level --- the mean of the
Monte-Carlo paths coincides with the analytic mean-reverting forecast, which is
the sanity check on the simulation. The dispersion of the grey paths shows
forecast uncertainty growing with horizon. The realized MA and EWMA lines
should fall within the simulated cone; systematic deviation (e.g. realized vol
persistently below the forecast mean) would indicate the forecast window was
calmer than the fitted regime implied. EWMA, being the most reactive of the
three, tracks the realized path most closely day to day, whereas the GARCH
forecast mean is the smoothest because it is an expectation, not a realization.

\bigskip
\noindent\textbf{Reproducing the figures.} From the repository root:
\texttt{pip install -r requirements.txt} then \texttt{python src/main.py}.
The script downloads and caches SPY prices, prints the GARCH fit, and writes
both figures into \texttt{figures/}; recompiling this document then embeds the
real-data plots automatically.

%----------------------------------------------------------------------
\appendix
\section{Python source code}
The complete, runnable package is hosted at the repository link on page~1.
Key modules are listed below.

\subsection*{\texttt{src/vol\_estimators.py}}
\lstinputlisting[language=Python]{src/vol_estimators.py}

\subsection*{\texttt{src/garch\_model.py}}
\lstinputlisting[language=Python]{src/garch_model.py}

\subsection*{\texttt{src/data.py}}
\lstinputlisting[language=Python]{src/data.py}

\subsection*{\texttt{src/main.py}}
\lstinputlisting[language=Python]{src/main.py}

\end{document}
